\section{What Makes a Decoder Work?}
\label{sec:repr}

% 证据：实验 0 结论 / 实验0_报告.md；排班_解码方案汇总.md（精确最优解的结构）

The six decoders of Section~\ref{sec:problem:decoders} differ widely in performance.
This section asks which property of the genotype--phenotype map explains the difference.
The answer motivates the design in Section~\ref{sec:method}.

\subsection{Measures}
\label{sec:repr:measures}

The phenotype is the dispatch sequence together with the types of the new tools; it fixes the roster.
For each instance and decoder we sample 1000 uniform random genotypes and compute:
\begin{itemize}
  \item \textbf{Neutral rate}: the probability that resampling one robot gene leaves the phenotype unchanged.
  It measures how many genotypes map to the same roster, \ie the redundancy of the decoder.
  \item \textbf{Initial supply}: the probability that a random genotype decodes to a roster with $\fres \le \LB + 1$.
  It measures how strongly the redundancy favors good rosters.
  \item \textbf{Cascade ratio}: after a single-gene change at shift $i$, the fraction of later shifts whose dispatch changes.
  It measures locality.
  \item \textbf{Reachability}: whether some genotype decodes to the exact optimal roster, which we obtain by dynamic programming on the four original scenarios.
\end{itemize}
Neutral rate and cascade ratio are measured on 100 random feasible genotypes and on the 30 best of the 1000 samples, with five single-gene changes each.
Performance is the hypervolume (HV) reached by NSGA-II \cite{deb2002nsga2} (population 30, about 465 evaluations).
For each instance we compute the Spearman correlation between a measure and HV across the decoders, and report the median over 16 instances
(the four original scenarios and 12 instances of the E, H and M families).

\subsection{Findings}
\label{sec:repr:findings}

\textbf{Redundancy helps when it favors good rosters.}
Table~\ref{tab:neutral} shows that each step from \dec{0} to \dec{5} makes the map more redundant.
On the original single-type scenario the neutral rate rises from 0.30 to 0.69; on the three multi-type scenarios it rises from about 0.1 to about 0.5.
The redundancy is not uniform: the initial supply rises with it.
Across instances, the median correlation with HV is $+0.91$ for the neutral rate (same sign on 12 of 12 instances)
and $+0.80$ for the initial supply (14 of 14).
This agrees with the theory of redundant representations \cite{rothlauf2006representations}:
non-uniform redundancy improves search when it over-represents good solutions.

\textbf{Locality is poor in every decoder and does not explain performance.}
A single-gene change alters 39--61\% of the later shifts, for all six decoders, with no trend from \dec{0} to \dec{5}.
The correlation of the cascade ratio with HV is only $+0.23$ (7 of 12).
In a sequential decoder, a changed decision changes the state, and every later gene is read in the new state.
All improvements so far came from better bias, not from better locality.

\textbf{Bias is useful but not robust.}
The capacity look-ahead of \dec{3}--\dec{5} replaces some \new actions by \reuse actions.
On two of the four scenarios this removes the optimal roster from the search space: the roster opens a fifth robot at shift 9, and the look-ahead never allows it.
Restoring reachability does not pay off.
When we add a gene band that forces \new, results get significantly worse on two scenarios and do not change on the other two.
The bias costs reachability, yet it is worth more than what it removes.
The same bias can still be wrong on other instances; Section~\ref{sec:exp-static} shows such a case.

\textbf{Well-balanced rosters are rarely sampled.}
Under every decoder, the probability that a random genotype gives $\fres = \LB$ and an imbalance within 1.5 times the optimum is at most 1.7\%,
and essentially zero on the three multi-type scenarios.
The remaining gap to the optimal front lies in this layer: at $\fres = \LB$, the best imbalance found is 1.2--2.3 times the optimum.

\begin{table}[t]
  \centering
  \caption{Neutral rate / initial supply $P(\fres \le \LB+1)$ of 1000 random genotypes on the four original scenarios.}
  \label{tab:neutral}
  \footnotesize
  \setlength{\tabcolsep}{1.6pt}
  \begin{tabular}{@{}lcccccc@{}}
    \toprule
    Scenario & \dec{0} & \dec{1} & \dec{2} & \dec{3} & \dec{4} & \dec{5} \\
    \midrule
    Single-type & 0.30/.000 & 0.42/.007 & 0.41/.007 & 0.61/.919 & 0.61/.929 & 0.69/.981 \\
    3 types, 720 h  & 0.10/.000 & 0.21/.000 & 0.22/.000 & 0.44/.020 & 0.46/.024 & 0.50/.032 \\
    3 types, 1440 h & 0.06/.000 & 0.13/.000 & 0.14/.000 & 0.38/.014 & 0.42/.008 & 0.44/.029 \\
    3 types, tight   & 0.12/.000 & 0.24/.000 & 0.19/.000 & 0.44/.020 & 0.47/.024 & 0.49/.032 \\
    \bottomrule
  \end{tabular}
\end{table}

\subsection{What Optimal Rosters Look Like}
\label{sec:repr:structure}

The exact optimal rosters have a simple structure.
About 90\% of their shifts follow one rule: send the available robot with the fewest duties.
The other 10\% deviate from it; in the single-type scenario, 10 of the 65 reuse decisions deviate.
The deviations sit at a few critical shifts; everywhere else the rule is right.

A per-shift encoding does not express this structure well.
Under \dec{5}, a uniform random genotype deviates from the ranking at many shifts, and each deviation changes the state seen by all later genes.
The rule itself also matters a great deal.
On the single-type scenario, following the fewest-duties rule without deviations needs 11--17 resources,
because the robots fall into maintenance together and leave gaps.
A rule that ranks robots by remaining budget reaches the lower bound of 9 without any search.

\subsection{Implications for Design}
\label{sec:repr:implications}

These findings lead to three requirements.
\begin{enumerate}
  \item The encoding should express ``a rule plus a few deviations'' directly. The genes should say where and how to deviate, and the rule should decide everything else.
  \item The strength of the bias should be a parameter. A fixed bias is useful but not robust, so the representation needs a knob that trades bias against freedom.
  \item The rule should be chosen for the bias it gives to search, not only for the quality of the single roster it builds.
\end{enumerate}
Section~\ref{sec:method} turns these requirements into a decoder.
Section~\ref{sec:theory} analyzes the knob.
